\subsection{An exact pure Weil realization of the actual family}
The actual algebra and module admit a finite realization by a
specified linear Frobenius. This retains the collision nilpotent
and both commuting actions. Its comparison with extension-dependent
quantum grading can then be tested directly.

\begin{finding}[\fid{F46}: actual End/Hom, coefficient geometry and pure Weil data]\label{f:cayleyweil}
Put
\[
 \delta=v+v^{-1},\quad D=\delta^2,\quad h=D-1,\quad
 z_1=\frac hD,\quad z_2=\frac{4h}{D^2},\quad
 \Delta=z_1-z_2=\frac{(D-1)(D-4)}{D^2},
\]
and use the finite coefficient localization
\[
 S_0=\Q(i)[v,v^{-1},\delta^{-1},h^{-1},(D-2)^{-1}],
 \qquad S=S_0[(1-\Delta)^{-1},(1+\Delta^2)^{-1}].
\]
Let $C=S[x]/(x(x-\Delta))$, with $x=z_1-z$ in the actual
presentation, and let $L_1=C/(x)$, $L_2=C/(x-\Delta)$.
Give the regular rank-two module $J=C$ the linear operator
\[
 F_J=\frac{1+ix}{1-ix}=1+s x,\qquad
 s=\frac{2i}{1-i\Delta},\qquad
 \lambda=\frac{1+i\Delta}{1-i\Delta}.
\]
For a scalar $a$, write $S_a$ for a rank-one module with
Frobenius $a$. Give $L_1,L_2$ Frobenius $1,\lambda$ and set
\[
 H=S_{-1}\oplus S_i\oplus J^{\oplus2},\qquad
 K=S_{-1}^{20}\oplus S_i^4\oplus L_1^{16}\oplus L_2^4.
\]
These have ranks $6,44$. Then
\[
 \End_F(H)=A=S\times S\times M_2(C),\qquad
 \operatorname{Hom}_F(H,K)=M,
\]
as the actual rank-three algebra and its rank-$64$ right tensor
module in \fid{F19}/\fid{F21}, with all named generators retained.
These identities commute with every coefficient base change.
Postcomposition also identifies
\[
 \End_F(K)\simeq\End_A(M),
\]
so the actual ordinary left operators from \fid{F43} lift to $K$.
The latter commutant has rank $688$ and special-fiber dimension
$816$ at $v=1$.

There are perfect coherent objects $\mathscr H,\mathscr K$ on
the smooth proper coefficient scheme $\mathbf P^1_S$ with those
same End/Hom groups and
\[
 \operatorname{Ext}^1_{\mathbf P^1_S}(\mathscr K,\mathscr K)
       =S^{688}\oplus(S/(\Delta))^{128}.
\]
Thus both crossed extension directions and the order-two collision
are retained, with additional self-Ext explicitly present.

For each prime power $q>1$, all evaluations $v=q^{\pm1/2}$ are
defined and give finite Weil representations pure of weight zero
under every complex embedding. At $v=1$, $J$ retains a nontrivial
Jordan block with eigenvalue one. These are chosen linear Weil
data; they are not scheme Frobenius on the coefficient geometry.
They do not satisfy extension-wise parameter retuning.
\end{finding}

\begin{proof}
\par\smallskip\noindent\textbf{The finite coefficient domain and cyclic operator.}\quad
The parameters $z_1,z_2$ and $\Delta$ belong to $S_0$. Both $1-i\Delta$ and
$1+i\Delta$ are units in $S$ because their product is
$1+\Delta^2$. For real $v\ne0$,
\[
 D\ge4,\qquad
 0\le\Delta=1-\frac5D+\frac4{D^2}<1.
\]
Thus every displayed denominator is nonzero at every such real
parameter, in particular at $q^{\pm1/2}$ and at $1$. This finite
localization admits the required evaluations; it does not
invert $\Delta$.

The relation $x^2=\Delta x$ gives
\[
 \frac{1+ix}{1-ix}=1+\frac{2i}{1-i\Delta}x,\qquad
 \lambda=1+s\Delta.
\]
The inverse is obtained by replacing $i$ by $-i$.
In the column basis $(1,x)$,
\begin{equation}\label{eq:cayleymatrix}
 F_J=\begin{pmatrix}1&0\\s&\lambda\end{pmatrix}.
\end{equation}
Since $s$ is a unit, multiplication by $x$ is
$s^{-1}(F_J-1)$. Therefore $S[F_J]=C$, and a linear map
commutes with $F_J$ exactly when it is $C$-linear.
The centralizer on the regular module is consequently
\[
 \End_{F_J}(J)=\End_C(C)=C.
\]
The argument works after every coefficient base change.
Equivalently $x=(F_J-1)/(i(F_J+1))$.
The vector $1$ is cyclic, since $(1,F_J1)$ has determinant $s$.
The characteristic polynomial is $(T-1)(T-\lambda)$.
At $v=1$ it remains cyclic and becomes
\[
 F_J=1+2ix,\qquad x^2=0,\quad x\ne0.
\]
Thus its repeated eigenvalue is a genuine Jordan block.

\par\smallskip\noindent\textbf{Exact algebra, module and both actions.}\quad
The scalar sectors are separated by units:
\[
 F_J+1=\frac2{1-ix},\qquad
 F_J-i=\frac{(1-i)(1-x)}{1-ix}.
\]
The second is invertible because the norm of $1-x$ is
$1-\Delta$, a unit of $S$. Also $-1-i$ is a unit.
Hence no intertwiners occur between different scalar sectors
or between those sectors and $J$. The same identities hold
on the branch quotients. The matrix sector of $\End_F(H)$ is
therefore $M_2(C)$, giving the asserted algebra.

An intertwiner from $J$ to a branch commutes with
$x=s^{-1}(F-1)$ and is $C$-linear. Hence
\[
 \operatorname{Hom}_F(J^{\oplus2},L_j)
    =\operatorname{Hom}_C(C^{\oplus2},L_j)
    =\operatorname{Row}(L_j).
\]
The scalar Hom groups contribute ranks $20,4$ and the branch
groups ranks $32,8$. In the column convention for $H$,
precomposition is right multiplication on these rows.
Thus this is precisely
\[
 M=\epsilon_+^{20}\oplus\epsilon_-^4
       \oplus\operatorname{Row}(L_1)^{16}
       \oplus\operatorname{Row}(L_2)^4.
\]
The chosen coordinate makes $L_1$ the multiplicity-$16$
branch $z=z_1$ and $L_2$ the multiplicity-$4$ branch $z=z_2$.
In particular the named normalized generator matrices are
\[
 \begin{pmatrix}1&0\\0&0\end{pmatrix},\qquad
 \begin{pmatrix}1-z&1\\z(1-z)&z\end{pmatrix},
 \qquad z=z_1-x,
\]
with the original generators obtained by multiplying by $D$,
and their scalar values $D,0$ as in \fid{F19}.
The certified regular lattice comparison of \fid{F21}
is defined over $S_0$, so this retains the actual action.
All arguments used unit identities and regular $C$-modules,
and therefore commute with arbitrary coefficient base change.
No diagonalization using $\Delta^{-1}$ was used.

Over the domain $S$, a map between the branch lines is
annihilated by $\lambda-1=s\Delta$ and hence is zero.
It follows that
\begin{equation}\label{eq:cayleycommutant}
 \End_F(K)=M_{20}(S)\times M_4(S)\times M_{16}(S)\times M_4(S).
\end{equation}
At $v=1$ the branch eigenvalues coincide and their multiplicities
combine into $20$, so this dimension increases from $688$ to
$400+16+400=816$. The base-change map is block diagonal
on those $20$ coordinates and has cokernel dimension $128$.
There is no such jump in $\End_F(H)$ or $\operatorname{Hom}_F(H,K)$,
because their source $J$ remains cyclic at collision.

Postcomposition gives an isomorphism
\[
 \End_F(K)\longrightarrow\End_A(\operatorname{Hom}_F(H,K)).
\]
For the scalar sectors this is immediate. On the matrix sector
put $K_C=L_1^{16}\oplus L_2^4$. The two diagonal matrix units
force an $M_2(C)$-commuting map on the row module $K_C^2$
to preserve its two coordinates; the two off-diagonal units
force the same $C$-linear map on each coordinate.
This is exactly postcomposition by a unique endomorphism of
$K_C$. The proof works over every base ring.
In the multiplicity bases of \fid{F43}, this is the identical
product-matrix action, so each actual ordinary left operator
lifts to $K$ and commutes with its chosen $F$.
Both commutants have dimension $816$ at collision, while
the additional $128$ maps need not lift over $S$.

\par\smallskip\noindent\textbf{Smooth proper coefficient geometry and extensions.}\quad
On $\mathbf P^1_S$ let $T$ denote the affine coordinate and put
\[
 Z=V((T-1)(T-\lambda)).
\]
It is a finite flat degree-two divisor, avoiding both infinity
and zero. The homomorphism
\[
 S[T,T^{-1}]/((T-1)(T-\lambda))
       \longrightarrow C,\qquad T\longmapsto1+s x
\]
is an isomorphism with inverse $x=s^{-1}(T-1)$, because its
relation becomes $s^2x(x-\Delta)=0$.
Write $O_a$ for the section sheaf at $T=a$, and define
\[
 \mathscr H=O_{-1}\oplus O_i\oplus O_Z^{\oplus2},\qquad
 \mathscr K=O_{-1}^{20}\oplus O_i^4
       \oplus O_1^{16}\oplus O_\lambda^4.
\]
Each section has its Cartier resolution by $O(-1),O$,
and $Z$ has the corresponding resolution by $O(-2),O$.
Thus these objects are perfect on the smooth proper
relative projective line. Their supports are finite over
affine $\Spec S$, so derived global sections are their
finite free modules in degree zero, namely $H,K$.
Computing sheaf Hom on the affine chart containing their
supports gives
\[
 \End_{\mathbf P^1_S}(\mathscr H)=A,\qquad
 \operatorname{Hom}_{\mathbf P^1_S}(\mathscr H,\mathscr K)=M.
\]
Multiplication by $T$ supplies the extra linear operators.

For the coherent Ext computation use the section resolution
\[
 0\longrightarrow O(-1)\xrightarrow{T-a}O
       \longrightarrow O_a\longrightarrow0.
\]
Hom into $O_b$ gives the two-term complex
$S\xrightarrow{b-a}S$. Its terms have no higher cohomology,
being supported on an affine section. Consequently
\[
 \operatorname{Hom}(O_a,O_b)=\operatorname{ann}_S(b-a),
 \qquad
 \operatorname{Ext}^1(O_a,O_b)=S/(b-a),
\]
and there are no higher Ext groups. For $a=b$ the two
groups are $S$. All distinct scalar-sector differences
are units; the sole nonunit cross difference is
$\lambda-1=s\Delta$. Both branch directions therefore give
$S/(\Delta)$, and
\begin{equation}\label{eq:cayleyext}
 \operatorname{Ext}^1_{\mathbf P^1_S}(\mathscr K,\mathscr K)
      =S^{688}\oplus(S/(\Delta))^{128}.
\end{equation}
The free rank is $20^2+4^2+16^2+4^2$; the torsion multiplicity
is $2\cdot16\cdot4$. At the actual order-two collision its
torsion length is $256$. This retains both crossed groups,
but the additional free self-Ext is not the self-Ext of the
original $C$-branch modules. No full derived equivalence
or extra Frobenius-equivariant derived Ext comparison is
asserted.

\par\smallskip\noindent\textbf{All-conjugate purity and continuous realizations.}\quad
Fix a prime power $q>1$ and evaluate $v=q^{1/2}$ or $q^{-1/2}$
in $\Q(i,\sqrt q)$. Then $D=q+2+q^{-1}$ and $\Delta$ is
rational with $0<\Delta<1$. The eigenvalues of $H,K$ lie in
\[
 \{-1,i,1,\lambda\},\qquad
 \lambda=\frac{1+i\Delta}{1-i\Delta}.
\]
Every complex embedding fixes $\Delta$ and sends $i$ to
$i$ or $-i$. Thus it sends $\lambda$ to $\lambda$ or
$\lambda^{-1}$, both of modulus one. The possible sign
change of $\sqrt q$ does not change $D$ or $\Delta$.
After an embedding into a finite extension of $\Q_\ell$,
these are finite Weil representations pure of weight zero
in the eigenvalue sense. No integrality assertion is
part of this purity conclusion.
For a general real algebraic parameter, total reality is a
sufficient additional hypothesis for the analogous all-conjugate
assertion; it is not asserted to be necessary.

At $v=1$ the same construction retains the nontrivial Jordan
block \eqref{eq:cayleymatrix} with eigenvalue one. It can be
viewed as a pure Weil object over any chosen finite field;
this does not posit a field of cardinality one.
Geometric restriction to a point forgets that arithmetic
Jordan data, so there is no contradiction with geometric
semisimplicity.

The ordinary internal Hom has dimension $6\cdot44=264$ and
Frobenius
\[
 \Phi(u)=F_KuF_H^{-1}.
\]
Its eigenvalues are ratios of the preceding unit-modulus
eigenvalues, so it too is pure of weight zero, with Jordan
blocks allowed. Its fixed subspace is exactly $M$, on which
$\Phi$ is the identity. Likewise $A$ is the fixed subalgebra
of the $36$-dimensional internal End object.

A finite Weil representation on a point is an invertible
linear operator, since the point Weil group is discrete
infinite cyclic. Extending it to a continuous representation
of the profinite Galois group requires its eigenvalues to be
$\ell$-adic units. This condition is necessary: compactness
of the image bounds all positive and negative powers.
It is also sufficient. After a finite coefficient extension,
choose Jordan form with unit eigenvalues. The matrix and
its inverse preserve a suitable integral lattice, so their
powers lie in a compact matrix group. Modulo each lattice
congruence subgroup the cyclic image is finite; these maps
extend compatibly from $\Z$ to $\widehat{\Z}$.
For each fixed $q$, the algebraic number $\lambda$ is a unit
outside finitely many finite places. Choosing
$\ell\nmid q$ outside that set therefore gives continuous
$\ell$-adic realizations retaining every displayed comparison.
No uniform choice for all $q$, or validity at every $\ell$,
is asserted.

\par\smallskip\noindent\textbf{A proved extension mismatch and an integrality boundary.}\quad
Write $\lambda(q)$ for the branch value at $v=q^{1/2}$.
For $q>1$,
\[
 \lambda(q)\ne1,\qquad
 \Delta(q^n)\longrightarrow1,\qquad
 \lambda(q^n)\longrightarrow i.
\]
If $\lambda(q)^n=\lambda(q^n)$ for every sufficiently
large $n$, the ratio of consecutive terms on the right
would tend to one, because their limit is nonzero.
The same ratio on the left is $\lambda(q)$, a contradiction.
Thus raising the fixed Frobenius to its $n$th power is
not retuning this named branch at $q^n$.

For a concrete arithmetic boundary, at $q=2$ one has
\[
 D=\frac92,\quad \Delta=\frac7{81},\quad
 \lambda=\frac{81+7i}{81-7i}
       =\frac{3256+567i}{3305},\quad
 \lambda+\lambda^{-1}=\frac{6512}{3305}\notin\Z[1/2].
\]
If $\lambda$ were integral over $\Z[1/2]$, its conjugate
and their rational sum would be integral, and that sum
would lie in the integrally closed ring $\Z[1/2]$.
This contradiction is unchanged by multiplying $\lambda$
by any integer power of $2$.
In particular it cannot be an eigenvalue in a Tate twist
of constant-coefficient cohomology of a smooth projective
variety over $\F_2$: the monic integral characteristic
polynomial in Deligne's theorem
\cite[Theorem~(1.6), p.~276; argument on pp.~276--277]{WeilI}
makes the untwisted eigenvalues algebraic integers, and
Tate twists introduce only powers of $2$ into denominators.
This excludes that extra comparison, not the constructed
pure Weil objects. For example $\lambda$ is a unit at
$\ell=3$ and gives a continuous realization there.
At $\ell=5$, the residues $i=\pm2$ make exactly one of
$81+7i$ and $81-7i$ divisible by $5$, so it is not a unit.
\end{proof}

The projective line above is over the characteristic-zero
parameter/coefficient base. It is not thereby a variety
over $\F_q$ with geometric Frobenius equal to multiplication
by $T$. The evaluation $v=\sqrt q$ occurs in a characteristic-zero
coefficient field and cannot simultaneously be that same
characteristic-$p$ base evaluation, where $q$ becomes zero
while $v$ must remain invertible. The coherent geometry and
the chosen linear Weil operator are separate data.

This realization closes the existence question for exact
finite pure coefficient objects and ordinary smooth proper
coefficient geometry retaining this rank-three action.
It does not identify the chosen operator with absolute
Frobenius on the center, the lift of \fid{F44}, or cyclotomic
Frobenius on TP. The unit-modulus spectrum does not supply
a positive self-adjoint structure, the designated VIII
globalization, or the \fid{F24} path-sign comparison.
The extension mismatch and integrality calculation above
are explicit limits of those additional interpretations.

% Separate paragraph for Part III.
\paragraph{\fid{F46} in Grothendieck's language.}
A finite flat divisor on $\mathbf P^1_S$ and its section
sheaves supply actual coherent coefficient objects with
Cartier perfect resolutions. Their endomorphisms, module
actions and crossed extensions preserve the specified
family and its collision order. Multiplication by a
coordinate gives separately chosen linear Weil data after
coefficient specialization; it is not the arithmetic
Frobenius of that coefficient scheme. This distinction
allows an exact pure End/Hom realization while keeping the
required grading, globalization and path comparisons as
additional conditions. No unpublished archive leaf is
claimed to contain the Cayley construction or these
nonstandard comparisons.

The affine module/sheaf dictionary is EGA~I, Theorem~1.3.7,
p.~86, Corollary~1.3.8, p.~87, and Theorem~1.4.1,
pp.~90--92. Perfect complexes use SGA~6, I, Definition~4.7,
pp.~121--122, and Example~4.8, p.~122. The equivariant-module
framework is SGA~3, I, 4.7.2--3, with cote~90, scan~17,
as the audited archive anchor. The eigenvalue and complex-weight
conventions are BBD~5.1.2 and~5.1.8. These locators supply
the framework; the displayed resolutions and matrix-unit
arguments prove this particular comparison.
